\newpage
\phantomsection
\label{prf:inverse-positive}
\LRAProofFor{thm:inverse-positive}

\begin{remark*}[Return]
\hyperref[thm:inverse-positive]{Return to Theorem}
\end{remark*}

\begin{theorem*}[Positive Elements Have Positive Inverses]
If $a>0$, then $a^{-1}>0$.
\end{theorem*}

\begin{proof}[Professional Standard Proof]
\LRAProofBodyStart
Let $F$ be an ordered field, and suppose that $a>0$.  Since $F$ is a field
and $a\neq 0$, the multiplicative inverse $a^{-1}$ exists.

We claim that $a^{-1}>0$.  Suppose, for contradiction, that
\[
a^{-1}<0.
\]
By adding $-a^{-1}$ to both sides, we obtain
\[
0<-a^{-1}.
\]
Now $0<a$ and $0<-a^{-1}$, so applying the order axiom for multiplication by
a positive element to the inequality
\[
0<-a^{-1}
\]
with multiplier $a$, we get
\[
0\cdot a<(-a^{-1})\cdot a.
\]
Thus
\[
0<a(-a^{-1}).
\]
Since
\[
a(-a^{-1})=-(aa^{-1})=-1,
\]
it follows that
\[
0<-1.
\]
Adding $1$ to both sides yields
\[
1<0.
\]
Now multiply the inequality $1<0$ by the positive element $a$.  Again by the
order axiom for multiplication by a positive element,
\[
1\cdot a<0\cdot a.
\]
Hence
\[
a<0,
\]
which contradicts the assumption that $a>0$.

Therefore our assumption $a^{-1}<0$ was false.  Since $a^{-1}\neq 0$, the
trichotomy law in the ordered field implies that
\[
a^{-1}>0.
\]
\end{proof}

\begin{proof}[Detailed Learning Proof]
\LRAProofBodyStart
We expand the professional proof by following the same sequence of justified moves.

Let $F$ be an ordered field, and suppose that $a>0$.  Since $F$ is a field
and $a\neq 0$, the multiplicative inverse $a^{-1}$ exists.

We claim that $a^{-1}>0$.  Suppose, for contradiction, that
\[
a^{-1}<0.
\]
By adding $-a^{-1}$ to both sides, we obtain
\[
0<-a^{-1}.
\]
Now $0<a$ and $0<-a^{-1}$, so applying the order axiom for multiplication by
a positive element to the inequality
\[
0<-a^{-1}
\]
with multiplier $a$, we get
\[
0\cdot a<(-a^{-1})\cdot a.
\]
Thus
\[
0<a(-a^{-1}).
\]
Since
\[
a(-a^{-1})=-(aa^{-1})=-1,
\]
it follows that
\[
0<-1.
\]
Adding $1$ to both sides yields
\[
1<0.
\]
Now multiply the inequality $1<0$ by the positive element $a$.  Again by the
order axiom for multiplication by a positive element,
\[
1\cdot a<0\cdot a.
\]
Hence
\[
a<0,
\]
which contradicts the assumption that $a>0$.

Therefore our assumption $a^{-1}<0$ was false.  Since $a^{-1}\neq 0$, the
trichotomy law in the ordered field implies that
\[
a^{-1}>0.
\]
\end{proof}

\begin{remark*}[Proof structure]
Assume $a>0$ and argue by contradiction.  If the inverse were negative, then
its negation would be positive.  Multiplying the positive inequality
$0<-a^{-1}$ by the positive element $a$ forces $0<-1$, and then translation
gives $1<0$.  Multiplying that inequality by the same positive element $a$
produces $a<0$, contradicting $a>0$.  So the inverse cannot be negative; by
trichotomy and nonzeroness of inverses, it must be positive.
\end{remark*}

\begin{dependencies}
\begin{itemize}
  \item \hyperref[def:abstract-ordered-field]{Definition~\ref*{def:abstract-ordered-field}}
\end{itemize}
\end{dependencies}

\clearpage
